\ifdefined\gkver\else\def\gkver{student}\fi
\documentclass[\gkver]{gaokaozhenti}
\begin{document}

\chapter{2001年上海}

\begin{enumerate}
\item
%% number: 1
%% paperName: 2001年上海高考第 1 题 5 分
%% typeId: 2
%% type: 多选题
%% chapter: 功和能
%% point: 动能定理
%% method: 无
%% score: 5
%% degree: 650
%% duplicateId: 0
%% body:
跳伞运动员在刚跳离飞机、其降落伞尚未打开的一段时间内，下列说法中正确的是\xzanswer{CD}

\fourchoices[answer=CD]
{空气阻力做正功}
{重力势能增加}
{动能增加}
{空气阻力做负功}

%% answer: CD
%% memo:
\memoanswer{
跳伞运动员跳离飞机，在尚未打开降落伞的这段时间内，运动员向下运动，重力对运动员做正功，重力势能减少，空气阻力对运动员做负功；由于重力大于空气阻力，运动员向下做加速运动，其动能增加，故 AB 错误 CD 正确。
}

\item
%% number: 2
%% paperName: 2001年上海高考第 2 题 5 分
%% typeId: 2
%% type: 多选题
%% chapter: 原子和原子核
%% point: 原子核式结构
%% method: 无
%% score: 5
%% degree: 850
%% duplicateId: 0
%% body:
卢瑟福原子核式结构理论的主要内容有\xzanswer{ACD}

\fourchoices[answer=ACD]
{原子的中心有个核，叫做原子核}
{原子的正电荷均匀分布在整个原子中}
{原子的全部正电荷和几乎全部质量都集中在原子核里}
{带负电的电子在核外绕着核旋转}

%% answer: ACD
%% memo:
\memoanswer{
由卢瑟福原子核结构理论即可判断本题答案。
}

\item
%% number: 3
%% paperName: 2001年上海高考第 3 题 5 分
%% typeId: 1
%% type: 单选题
%% chapter: 电场
%% point: 电势能电势差电场力做功
%% method: 无
%% score: 5
%% degree: 650
%% duplicateId: 0
%% body:
A、B 两点各放有电量为 $+Q$ 和 $+2Q$ 的点电荷，A、B、C、D 四点在同一直线上，且 $AC=CD=DB$。将一正电荷从 C 点沿直线移到 D 点，则\xzanswer{B}

\onepicture[align=right, w=9cm]{\includesvg{figs/03}}

\fourchoices[answer=B]
{电场力一直做正功}
{电场力先做正功再做负功}
{电场力一直做负功}
{电场力先做负功再做正功}

%% answer: B
%% memo:
\memoanswer{
由于 A、B 两点均放有带正电的点电荷，在 AB 连线上有一点合场强必定为零，设该点为 P，且距 A 点距离为 $x$，AB 两点间的距离为 $L$，由点电荷场强公式和场的叠加原理有 $E_{P}=k\frac{Q}{x^{2}}-k\frac{2Q}{(L-x)^{2}}=0$，解得 $x=(\sqrt{2}-1)L$，可见 P 点在 CD 之间，P 点左侧合场强方向水平向右，P 点右侧合场强方向水平向左，所以将一正电荷从 C 点沿直线移到 D 点的过程中，电场力先做正功再做负功，故 B 正确。
}

\item
%% number: 4
%% paperName: 2001年上海高考第 4 题 5 分
%% typeId: 2
%% type: 多选题
%% chapter: 曲线运动
%% point: 天体运动
%% method: 无
%% score: 5
%% degree: 450
%% duplicateId: 0
%% body:
组成星球的物质是靠引力吸引在一起的，这样的星球有一个最大的自转速率，如果超过了该速率，星球的万有引力将不足以维持其赤道附近的物体做圆周运动。由此能得到半径为 $R$、密度为 $\rho$、质量为 $M$ 且均匀分布的星球的最小自转周期 $T$。下列表达式中正确的是\xzanswer{AD}

\fourchoices[answer=AD]
{$T=2\pi\sqrt{\frac{R^{3}}{GM}}$}
{$T=2\pi\sqrt{\frac{3R^{3}}{GM}}$}
{$T=\sqrt{\frac{\pi}{G\rho}}$}
{$T=\sqrt{\frac{3\pi}{G\rho}}$}

%% answer: AD
%% memo:
\memoanswer{
该星球对赤道附近物体的万有引力提供向心力，有 $G\frac{Mm}{R^{2}}=m\frac{4\pi^{2}}{T^{2}}R$，得 $T=2\pi\sqrt{\frac{R^{3}}{GM}}$，故 A 正确；又 $M=\rho\cdot\frac{4}{3}\pi R^{3}$，代入得 $T=\sqrt{\frac{3\pi}{\rho G}}$，故 D 正确。
}

\item
%% number: 5
%% paperName: 2001年上海高考第 5 题 5 分
%% typeId: 2
%% type: 多选题
%% chapter: 电磁感应
%% point: 电磁感应中的变加速
%% method: 无
%% score: 5
%% degree: 450
%% duplicateId: 0
%% body:
如图所示，有两根和竖直方向成 $\alpha$ 角的光滑平行的金属轨道，上端接有可变电阻 $R$，下端足够长，空间有垂直于轨道平面的匀强磁场，磁感强度为 $B$。一根质量为 $m$ 的金属杆从轨道上由静止滑下。经过足够长的时间后，金属杆的速度会趋近于一个最大速度 $v_{max}$，则\xzanswer{BC}

\onepicture[align=right, w=5cm]{\includesvg{figs/05}}

\fourchoices[answer=BC]
{如果 $B$ 增大，$v_{max}$ 将变大}
{如果 $\alpha$ 变小，$v_{max}$ 将变大}
{如果 $R$ 变大，$v_{max}$ 将变大}
{如果 $m$ 变小，$v_{max}$ 将变大}

%% answer: BC
%% memo:
\memoanswer{
由题意知，金属杆最终做匀速运动，由受力分析知，金属杆受重力、轨道的支持力和安培力的作用，安培力 $F=\frac{B^{2}L^{2}v_{m}}{R}$，由平衡条件得 $mg\sin\alpha=\frac{B^{2}L^{2}v_{m}}{R}$，解得 $v_{m}=\frac{mg\sin\alpha\cdot R}{B^{2}L^{2}}$，由此式可知，$B$ 增大，$v_{m}$ 减小；$\alpha$ 增大，$v_{m}$ 增大；$R$ 变大，$v_{m}$ 变大；$m$ 变小，$v_{m}$ 变小，故 AD 错误，BC 正确。
}

\item
%% number: 6
%% paperName: 2001年上海高考第 6 题 6 分
%% typeId: 2
%% type: 多选题
%% chapter: 电磁感应
%% point: 楞次定律
%% method: 无
%% score: 6
%% degree: 450
%% duplicateId: 0
%% body:
如图所示是一种延时开关，当 $S_{1}$ 闭合时，电磁铁 F 将衔铁 D 吸下，C 线路接通。当 $S_{1}$ 断开时，由于电磁感应作用，D 将延迟一段时间才被释放。则\xzanswer{BC}

\onepicture[align=right, w=5cm]{\includesvg{figs/06}}

\fourchoices[answer=BC]
{由于 A 线圈的电磁感应作用，才产生延时释放 D 的作用}
{由于 B 线圈的电磁感应作用，才产生延时释放 D 的作用}
{如果断开 B 线圈的电键 $S_{2}$，无延时作用}
{如果断开 B 线圈的电键 $S_{2}$，延时将变长}

%% answer: BC
%% memo:
\memoanswer{
延时开关的工作原理是：当断开 $S_{1}$，使 A 线圈中电流变小并消失时，铁芯中的磁通量减小，若 B 线圈闭合则在其中引起感应电流，由楞次定律得，感应电流的磁场阻碍原磁场的减小，因此就使铁芯中磁场减弱得慢些，因此才产生延时释放 D 的作用，可见是由于 B 线圈的电磁感应作用，起到了延时作用，故 BC 正确。
}

\item
%% number: 7
%% paperName: 2001年上海高考第 7 题 5 分
%% typeId: 1
%% type: 单选题
%% chapter: 电路
%% point: 故障分析
%% method: 无
%% score: 5
%% degree: 450
%% duplicateId: 0
%% body:
如图所示的电路中，闭合电键，灯 $L_{1}$、$L_{2}$ 正常发光，由于电路出现故障，突然发现灯 $L_{1}$ 变亮，灯 $L_{2}$ 变暗，电流表的读数变小，根据分析，发生的故障可能是\xzanswer{A}

\onepicture[align=right, w=5cm]{\includesvg{figs/07}}

\fourchoices[answer=A]
{$R_{1}$ 断路}
{$R_{2}$ 断路}
{$R_{3}$ 短路}
{$R_{4}$ 短路}

%% answer: A
%% memo:
\memoanswer{
等效电路如图所示。

\onepicture[align=center, w=7cm]{figs/07_an.jpg}

若 $R_{1}$ 断路，总外电阻变大，总电流减小，路端电压变大，$L_{1}$ 两端电压变大，$L_{1}$ 变亮；bc 部分电路结构没变，电流仍按原比例分配，总电流减小，通过 $L_{2}$、电流表的电流都减小，故 A 正确；若 $R_{2}$ 断路，总外电阻变大，总电流减小，ab 部分电路结构没变，电流仍按原比例分配，$L_{1}$ 中电流减小，与题意相矛盾，故 B 错误；若 $R_{3}$ 短路或 $R_{4}$ 短路，总外电阻减小，总电流增大，电流表中电流变大，与题意相矛盾，故 CD 错误。
}

\item
%% number: 8
%% paperName: 2001年上海高考第 8 题 5 分
%% typeId: 2
%% type: 多选题
%% chapter: 运动和力
%% point: 变加速直线运动
%% method: 无
%% score: 5
%% degree: 450
%% duplicateId: 0
%% body:
一升降机在箱底装有若干个弹簧，设在某次事故中，升降机吊索在空中断裂，忽略摩擦力，则升降机在从弹簧下端触地后直到最低点的一段运动过程中，\xzanswer{CD}

\onepicture[align=right, w=3cm]{\includesvg{figs/08}}

\fourchoices[answer=CD]
{升降机的速度不断减小}
{升降机的加速度不断变大}
{先是弹力做的负功小于重力做的正功，然后是弹力做的负功大于重力做的正功}
{到最低点时，升降机加速度的值一定大于重力加速度的值}

%% answer: CD
%% memo:
\memoanswer{
由胡克定律知，弹簧对升降机的弹力 $F$ 随弹簧的形变量 $x$ 的增大而增大。①在重力大于弹力过程中，由牛顿第二定律得 $mg-kx=ma$，$x$ 增大，弹力变大，合外力变小，加速度变小，升降机做加速度逐渐变小的加速运动，弹力做的负功小于重力做的正功；②当重力等于弹力时，合外力等于零，加速度为零，此时升降机速度达到最大值；③在重力小于弹力过程中，有 $kx-mg=ma$，$x$ 增大，弹力增大，合外力增大，加速度增大，升降机做加速度逐渐变大的减速运动，弹力做的负功大于重力做的正功，故 AB 错误 CD 正确。
}

\item
%% number: 9
%% paperName: 2001年上海高考第 9 题 4 分
%% typeId: 3
%% type: 填空题
%% chapter: 其他
%% point: 物理学史
%% method: 无
%% score: 4
%% degree: 850
%% duplicateId: 0
%% body:
请将右面三位科学家的姓名按历史年代先后顺序排列：\tkanswer[1.4cm]{伽利略}、\tkanswer[1.4cm]{牛顿}、\tkanswer[1.4cm]{爱因斯坦}。任选其中二位科学家，简要写出他们在物理学上的主要贡献各一项：\tkanswer[3.4cm]{伽利略：望远镜的早期发明，将实验方法引进物理学等}，\tkanswer[3.4cm]{牛顿：发现运动定律，万有引力定律等}。

\threepicture[align=center, showprefix=false, h=3.2cm, capA=牛顿, capB=伽利略, capC=爱因斯坦]
{figs/09a.jpg}
{figs/09b.jpg}
{figs/09c.jpg}

%% answer:
%% memo:
\memoanswer{
伽利略，牛顿，爱因斯坦。

伽利略：望远镜的早期发明，将实验方法引进物理学等；

牛顿：发现运动定律，万有引力定律等；

爱因斯坦：光电效应，相对论等。
}

\item
%% number: 10
%% paperName: 2001年上海高考第 10 题 4 分
%% typeId: 3
%% type: 填空题
%% chapter: 光学
%% point: 光的衍射
%% method: 无
%% score: 4
%% degree: 650
%% duplicateId: 0
%% body:
（a）、（b）两幅图是由单色光分别射到圆孔而形成的图象，其中图\ref{2001上海10a}是光的\tkanswer[1.4cm]{衍射}（填干涉或衍射）图象。由此可以判断出图\ref{2001上海10a}所对应的圆孔的孔径\tkanswer[1.4cm]{小于}（填大于或小于）图\ref{2001上海10b}所对应的圆孔的孔径。

\twopicture[align=right, w=3.2cm, labelA=2001上海10a, labelB=2001上海10b]
{figs/10a.png}
{figs/10b.png}

%% answer:
%% memo:
\memoanswer{
图\ref{2001上海10a}是光的衍射图样，图\ref{2001上海10b}是小孔成像，由衍射条件知，图\ref{2001上海10a}所对应的圆孔的孔径小于图\ref{2001上海10b}。
}

\item
%% number: 11
%% paperName: 2001年上海高考第 11 题 4 分
%% typeId: 3
%% type: 填空题
%% chapter: 电场
%% point: 带电粒子的偏转
%% method: 无
%% score: 4
%% degree: 650
%% duplicateId: 0
%% body:
一束质量为 $m$、电量为 $q$ 的带电粒子以平行于两极板的速度 $v_{0}$ 进入匀强电场，如图所示，如果两极板间电压为 $U$，两极板间的距离为 $d$，板长为 $L$，设粒子束不会击中极板，则粒子从进入电场到飞出极板时电势能的变化量为\tkanswer[2.6cm]{$\frac{q^{2}U^{2}L^{2}}{2md^{2}v_{0}^{2}}$}（粒子的重力忽略不计）

\onepicture[align=right, w=5.5cm]{\includesvg{figs/11}}

%% answer:
%% memo:
\memoanswer{
带电粒子以 $v_{0}$ 速度进入匀强电场，在 $v_{0}$ 方向上做匀速直线运动；在垂直于极板方向上受电场力的作用，做初速度为零的匀加速运动，有加速度 $a=\frac{qU}{dm}$，时间 $t=\frac{L}{v_{0}}$，末垂直于极板方向的速度 $v=at=\frac{qUL}{dmv_{0}}$，电场力对粒子做正功，粒子的动能增加，由能量守恒可知，粒子电势能的变化量 $\Delta E_{p}$ 等于粒子动能的增量 $\Delta E_{k}$，则有 $\Delta E_{p}=\Delta E_{k}$，$\Delta E_{k}=\frac{1}{2}mv^{2}=\frac{q^{2}U^{2}L^{2}}{2md^{2}v_{0}^{2}}$，所以粒子从进入电场到飞出极板时，电势能的变化量为 $\frac{q^{2}U^{2}L^{2}}{2md^{2}v_{0}^{2}}$。
}

\item
%% number: 12
%% paperName: 2001年上海高考第 12 题 4 分
%% typeId: 3
%% type: 填空题
%% chapter: 机械波
%% point: 波长频率波速
%% method: 无
%% score: 4
%% degree: 450
%% duplicateId: 0
%% body:
如图所示，有四列简谐波同时沿 $x$ 轴正方向传播，波速分别是 $v$、$2v$、$3v$ 和 $4v$，a、b 是 $x$ 轴上所给定的两点，且 $ab=l$。在 $t$ 时刻 a、b 两点间四列波的波形分别如图所示，则由该时刻起 a 点出现波峰的先后顺序依次是图\tkanswer[1.6cm]{BDCA}；频率由高到低的先后顺序依次是图\tkanswer[1.6cm]{DBCA}。

\fourpicture[align=center, w=3.2cm, numstyle=letter, labelprefix={}, numA=A, numB=B, numC=C, numD=D]
{figs/12a.png}
{figs/12b.png}
{figs/12c.png}
{figs/12d.png}

%% answer:
%% memo:
\memoanswer{
从波形图知 $\lambda_{A}=1,\lambda_{B}=\frac{1}{2},\lambda_{C}=2,\lambda_{D}=\frac{2}{3}$；由 $v=\lambda f$ 知 $f=\frac{v}{\lambda}$，有 $f_{A}=1,f_{B}=4,f_{C}=\frac{3}{2},f_{D}=6$，所以频率由高到低的先后顺序依次是图 DBCA。

把波形图向左延长，如图所示，由图象可以看出，由该时刻起 a 点出现波峰，最少要传播的距离为 $\frac{\lambda_{A}}{4}$，$\frac{\lambda_{B}}{4}$，$\frac{\lambda_{C}}{4}$，$\frac{3\lambda_{D}}{4}$，也就是该时刻起 a 点出现波峰，最少要传播的时间是 $\frac{T_{A}}{4}$，$\frac{T_{B}}{4}$，$\frac{T_{C}}{4}$，$\frac{3T_{D}}{4}$，因 $T=\frac{1}{f}$，所以 $T_{A}=1,T_{B}=\frac{1}{4},T_{C}=\frac{2}{3},T_{D}=\frac{1}{6}$，所以该时刻起 a 点出现波峰，最少要传播的时间是 $\frac{1}{4},\frac{1}{16},\frac{1}{6},\frac{1}{8}$，所以由该时刻起 a 点出现波峰的先后顺序依次是图 BDCA。

\fourpicture[align=center, showlabel=false, w=3.2cm]
{figs/12_an1.jpg}
{figs/12_an2.jpg}
{figs/12_an3.jpg}
{figs/12_an4.jpg}
}

\item
%% number: 13
%% paperName: 2001年上海高考第 13 题 4 分
%% typeId: 3
%% type: 填空题
%% chapter: 直线运动
%% point: 匀速直线运动
%% method: 无
%% score: 4
%% degree: 350
%% duplicateId: 0
%% body:
图\ref{2001上海13a}是在高速公路上用超声波测速仪测量车速的示意图，测速仪发出并接收超声波脉冲信号。根据发出和接收到的信号间的时间差，测出被测物体的速度。图\ref{2001上海13b}中 $p_{1}$、$p_{2}$ 是测速仪发出的超声波信号，$n_{1}$、$n_{2}$ 分别是 $p_{1}$、$p_{2}$ 由汽车反射回来的信号。设测速仪匀速扫描，$p_{1}$、$p_{2}$ 之间的时间间隔 $\Delta t=1.0\Us$，超声波在空气中传播的速度是 $v=340\Ums$，若汽车是匀速运动的，则根据图可知，汽车在接收到 $p_{1}$、$p_{2}$ 两个信号之间的时间内前进的距离是\tkanswer[1.2cm]{17}，汽车的速度是\tkanswer[1.2cm]{17.9}$\Ums$。

\twopicture[align=center, w=9cm, labelA=2001上海13a, labelB=2001上海13b]
{figs/13a.png}
{figs/13b.png}

%% answer:
%% memo:
\memoanswer{
（1）作运动情景分析示意图，图中，$P_{1}N_{1}$ 中间时刻是汽车接收 $P_{1}$ 时刻发出的超声波信号的时刻，此位置 A 是汽车的初位置；$P_{2}N_{2}$ 中间时刻是汽车接收 $P_{2}$ 时刻发出的超声波信号的时刻，此位置 B 是汽车的末位置，要求的汽车前进的距离就是 A、B 之间的距离。

\onepicture[align=center, w=8cm]{figs/13_an.jpg}

（2）求解

第 1 步，先读出图\ref{2001上海13b}中各时刻对应的读数

\begin{center}
\begin{tabular}{|c|c|c|c|c|c|c|}
\hline
点 & $P_{1}$ & $N_{1}$ & $P_{1}N_{1}$ 中点 & $P_{2}$ & $N_{2}$ & $P_{2}N_{2}$ 中点 \\
\hline
读数（格） & 0.5 & 1.7 & 1.1 & 3.5 & 4.4 & 3.95 \\
\hline
对应的时刻（$\Us$） & 0.17 & 0.57 & 0.37 & 1.17 & 1.47 & 1.32 \\
\hline
\end{tabular}
\end{center}

第 2 步，再算各点对应的时间，算法如下：先算每 1 大格对应的时间，$t_{1}=\frac{\Delta t}{P_{2}-P_{1}}=\frac{1.0}{3.5-0.5}=\frac{1}{3}$，再算各点对应的时间，$t=\frac{1}{3}\times$ 格数，填入上表的第 3 行。

第 3 步，算出 A 与 O 的距离

$s_{1}=340\times(0.37-0.17)\Um=68\Um$，

第 4 步，算出 B 与 O 的距离

$s_{2}=340\times(1.32-1.17)\Um=51\Um$，

第 5 步，算出 A 与 B 的距离

$\Delta s=s_{1}-s_{2}=68\Um-51\Um=17\Um$，

这就是汽车在接收 $P_{1}$、$P_{2}$ 两个时刻发出的超声波信号中间前进的距离。

第 6 步，算出汽车从 A 到 B 用的时间为

$\Delta t=1.32\Us-0.37\Us=0.95\Us$，

第 7 步，算出汽车的速度为

$v=\frac{\Delta s}{\Delta t}=\frac{17}{0.95}=17.9\Ums=64\Ukmh$。
}

\item
%% number: 14
%% paperName: 2001年上海高考第 14 题 5 分
%% typeId: 4
%% type: 实验题
%% chapter: 光学
%% point: 光电效应
%% method: 无
%% score: 5
%% degree: 650
%% duplicateId: 0
%% body:
光电效应实验的装置如图所示，则下面说法中正确的是\xzanswer{AD}

\onepicture[align=right, w=5cm]{figs/14.png}

\fourchoices[answer=AD]
{用紫外光照射锌板，验电器指针会发生偏转}
{用红色光照射锌板，验电器指针会发生偏转}
{锌板带的是负电荷}
{使验电器指针发生偏转的是正电荷}

%% answer: AD
\jdanswer{AD}
%% memo:
\memoanswer{
将擦得很亮的锌板连接验电器，用弧光灯照射锌板（弧光灯发出紫外线），验电器指针张开一个角度，说明锌板带了电，进一步研究表明锌板带正电，这说明在紫外线的照射下，锌板中有一部分自由电子从表面飞出来，锌板中缺少电子，于是带正电，故 AD 正确。
}

\item
%% number: 15
%% paperName: 2001年上海高考第 15 题 5 分
%% typeId: 4
%% type: 实验题
%% chapter: 气体性质
%% point: 验证玻意耳定律
%% method: 无
%% score: 5
%% degree: 650
%% duplicateId: 0
%% body:
某同学用同一个注射器做了两次验证玻意耳定律的实验，操作完全正确。根据实验数据却在 $p-V$ 图上画出了两条不同双曲线。造成这种情况的可能原因是\xzanswer{AB}

\onepicture[align=right, w=4.5cm]{\includesvg{figs/15}}

\fourchoices[answer=AB]
{两次实验中空气质量不同}
{两次实验中温度不同}
{两次实验中保持空气质量、温度相同，但所取的气体压的数据不同}
{两次实验中保持空气质量、温度相同，但所取的气体体积的数据不同}

%% answer: AB
\jdanswer{AB}
%% memo:
\memoanswer{
由理想气体状态方程 $\frac{PV}{T}=C$ 得 $PV=CT$，若 $PV$ 乘积一定，则 $p-V$ 图象为双曲线，且乘积不同，双曲线不同，故题中可能是温度不同，也可能是常数 $C$ 不同，而常数 $C$ 由气体质量决定，即也可能是气体质量不同，故 AB 正确 CD 错误。
}

\item
%% number: 16
%% paperName: 2001年上海高考第 16 题 6 分
%% typeId: 4
%% type: 实验题
%% chapter: 电路
%% point: 测电动势与内阻
%% method: 无
%% score: 6
%% degree: 450
%% duplicateId: 0
%% body:
要求测量由 2 节干电池串联而成的电池组的电动势 $E$ 和内阻 $r$（约几欧），提供下列器材：电压表 V（量程 $3\UV$，内阻 $1\UkO$）、电压表 $V_{2}$（量程 $15\UV$，内阻 $2\UkO$）、电阻箱（$0\sim9999\UO$）、电键、导线若干。

某同学用量程为 $15\UV$ 的电压表连接成如图所示的电路，实验步骤如下：
\begin{enumerate}
	\item
	合上电键 S，将电阻箱 $R$ 阻值调到 $R_{1}=10\UO$，读得电压表的读数为 $U_{1}$。
	\item
	将电阻箱 $R$ 阻值调到 $R_{2}=20\UO$，读得电压表的读数为 $U_{2}$。
\end{enumerate}

由方程组 $U_{1}=E-\frac{U_{1}}{R_{1}}r$，$U_{2}=E-\frac{U_{2}}{R_{2}}r$，解出 $E$、$r$。

为了减少实验误差，上述实验在选择器材和实验步骤中，应做哪些改进？

\onepicture[align=right, w=5cm]{\includesvg{figs/16}}

%% answer:
\jdanswer{
\begin{enumerate}
	\item
	应选用量程为 $3\UV$ 的电压表。
	\item
	改变电阻箱阻值 $R$，读取若干个 $U$ 的值，由 $I=\frac{U}{R}$ 计算出电流的值，然后作出 $U-I$ 图线，得到 $E$、$r$。
\end{enumerate}
}
%% memo:
\memoanswer{
因为 2 节干电池串联而成的电池组的电动势约 $3\UV$，所以应选用量程为 $3\UV$ 的电压表。为了减小偶然误差，要采用测量多组数据求平均值的方法，所以改变电阻箱阻值 $R$，读取若干个 $U$ 的值，由 $I=\frac{U}{R}$ 计算出电流的值，然后作出 $U-I$ 图线，得到 $E$、$r$。

也可以改变电阻箱阻值 $R$，读取若干个 $R$ 和 $U$ 的值，由 $E=U+\frac{U}{R}r$ 得到 $\frac{1}{U}=\frac{1}{E}+\frac{r}{E}R$，然后作出 $\frac{1}{U}-R$ 图线，得到 $E$、$r$。（截距的倒数等于电动势 $E$，截距的倒数与斜率的乘积等于内阻 $r$）。
}

\item
%% number: 17
%% paperName: 2001年上海高考第 17 题 7 分
%% typeId: 4
%% type: 实验题
%% chapter: 直线运动
%% point: 打点计时器公式
%% method: 无
%% score: 7
%% degree: 450
%% duplicateId: 0
%% body:
利用打点计时器研究一个约 $1.4\Um$ 高的商店卷帘窗的运动。将纸带粘在卷帘底部，纸带通过打点计时器随帘在竖直面内向上运动。打印后的纸带如下图所示，数据如表格所示。纸带中 AB、BC、CD……每两点之间的时间间隔为 $0.10\Us$，根据各间距的长度，可计算出卷帘窗在各间距内的平均速度 $v_{\text{平均}}$。可以将 $v_{\text{平均}}$ 近似地作为该间距中间时刻的即时速度 $v$。

\begin{center}
\begin{tabular}{|c|c|}
\hline
\multicolumn{2}{|c|}{卷帘运动的数据} \\
\hline
间隔 & 间距（$\Ucm$） \\
\hline
AB & 5.0 \\ \hline
BC & 10.0 \\ \hline
CD & 15.0 \\ \hline
DE & 20.0 \\ \hline
EF & 20.0 \\ \hline
FG & 20.0 \\ \hline
GH & 20.0 \\ \hline
HI & 17.0 \\ \hline
IJ & 8.0 \\ \hline
JK & 4.0 \\ \hline
\end{tabular}
\end{center}

\begin{enumerate}
	\item
	请根据所提供的纸带和数据，绘出卷帘窗运动的 $v-t$ 图像。
	\item
	AD 段的加速度为\tkanswer[1.2cm]{5}$\Umsq$，AK 段的平均速度为\tkanswer[1.2cm]{1.39}$\Ums$。
\end{enumerate}

\onepicture[align=right, h=7cm]{\includesvg{figs/17a}}

\onepicture[align=right, w=7cm]{\includesvg{figs/17b}}

%% answer:
\jdanswer{
\begin{enumerate}
	\item
	如图所示
	\item
	$5\Umsq$，$1.39\Ums$
\end{enumerate}
}
%% memo:
\memoanswer{
（1）求各间距内的平均速度 $\bar{v}=\frac{s_{i}}{\Delta t}$（$s_{i}$ 表示每两计时点间的距离，$\Delta t$ 表示每两计时点间的时间间隔），由已知条件可确定 $0.5\Us,1.5\Us,2.5\Us,\dots$ 各时刻的计时速度，计算数据列入下面表格中，绘出 $v-t$ 图，参见答案。

\begin{center}
\begin{tabular}{|c|c|c|c|c|c|c|c|c|c|c|}
\hline
时刻（$\Us$） & 0.5 & 1.5 & 2.5 & 3.5 & 4.5 & 5.5 & 6.5 & 7.5 & 8.5 & 9.5 \\
\hline
速度（$\Ums$） & 0.50 & 1.00 & 1.50 & 2.00 & 2.00 & 2.00 & 2.00 & 1.70 & 0.80 & 0.40 \\
\hline
\end{tabular}
\end{center}

\onepicture[align=center, w=7cm]{\includesvg{figs/17_an}}

（2）由 $v-t$ 图可见 AD 段卷帘窗做匀加速运动，由公式 $\bar{v}_{DE}-\bar{v}_{AB}=a\cdot3\Delta t$，或逐差法 $a=\frac{(DE+CD)-(AB+BC)}{(2\Delta t)^{2}}$ 均可求得 $a=5\Umsq$。AK 段的平均速度 $v_{AK}=\frac{AK}{10\Delta t}=\frac{139\times10^{-2}}{10\times0.1}\Ums=1.39\Ums$。
}

\item
%% number: 18
%% paperName: 2001年上海高考第 18 题 7 分
%% typeId: 4
%% type: 实验题
%% chapter: 电路
%% point: 分压与分流
%% method: 无
%% score: 7
%% degree: 450
%% duplicateId: 0
%% body:
某学生为了测量一物体的质量，找到一个力电转换器，该转换器的输出电压正比于受压面的压力（比例系数为 $k$），如图所示。测量时先调节输入端的电压，使转换器空载时的输出电压为 0；而后在其受压面上放一物体，即可测得与物体的质量成正比的输出电压 $U$。现有下列器材：力电转换器、质量为 $m_{0}$ 的砝码、电压表、滑动变阻器、干电池各一个、电键及导线若干、待测物体（可置于力电转换器的受压面上）。请完成对该物体质量的测量。
\begin{enumerate}
	\item
	设计一个电路，要求力电转换器的输入电压可调，并且使电压的调节范围尽可能大，在方框中画出完整的测量电路图。
	\item
	简要说明测量步骤，求出比例系数 $k$，并测出待测物体的质量 $m$。
	\item
	请设想实验中可能会出现的一个问题。
\end{enumerate}

\onepicture[align=right, w=5.5cm]{\includesvg{figs/18}}

%% answer:
\jdanswer{
\begin{enumerate}
	\item
	设计的电路图如图所示。
	\item
	测量步骤与结果：①调节滑动变阻器，使转换器的输出电压为零；②将砝码放在转换器上，记下输出电压 $U_{0}$；③将待测物放在转换器上，记下输出电压 $U_{1}$。由 $U_{0}=km_{0}g$，得 $k=\frac{U_{0}}{m_{0}g}$；测得 $U=kmg$，所以 $m=\frac{U}{U_{0}}m_{0}$。
	\item
	①因电源电压不够而输出电压调不到零；②待测物体质量超出转换器量程。
\end{enumerate}
}
%% memo:
\memoanswer{
（1）设计的电路图如图所示。

\onepicture[align=center, w=6cm]{\includesvg{figs/18_an}}

（2）测量步骤与结果：

①调节滑动变阻器，使转换器的输出电压为零；

②将砝码放在转换器上，记下输出电压 $U_{0}$；

③将待测物放在转换器上，记下输出电压 $U_{1}$。

由 $U_{0}=km_{0}g$，得 $k=\frac{U_{0}}{m_{0}g}$；测得 $U=kmg$，所以 $m=\frac{m_{0}U}{U_{0}}$。

（3）①因电源电压不够而输出电压调不到零；②待测物体质量超出转换器量程。

本题考查包含传感器的实验电路的设计、测试和测量及设计实验能力和创新能力，要求力电转换器的输入电压可调，并且使电压的调节范围尽可能大，所以把滑动变阻器接在输入端并且接成分压电路。
}

\item
%% number: 19
%% paperName: 2001年上海高考第 19 题 10 分
%% typeId: 6
%% type: 计算题
%% chapter: 功和能
%% point: 能量守恒定律
%% method: 建立模型
%% score: 10
%% degree: 450
%% duplicateId: 0
%% body:
1791 年，米被定义为：在经过巴黎的子午线上，取从赤道到北极长度的一千万分之一。
\begin{enumerate}
	\item
	请由此估算地球的半径 $R$。（答案保留二位有效数字）
	\item
	太阳与地球的距离为 $1.5\times10^{11}\Um$，太阳光以平行光束入射到地面。地球表面 $\frac{2}{3}$ 的面积被水面所覆盖，太阳在一年中辐射到地球表面水面部分的总能量 $W$ 约为 $1.87\times10^{24}\UJ$。设水面对太阳辐射的平均反射率为 $7\%$，而且将吸收到的 $35\%$ 能量重新辐射出去。太阳辐射可将水面的水蒸发（设在常温、常压下蒸发 $1\Ukg$ 水需要 $2.2\times10^{6}\UJ$ 的能量），而后凝结成雨滴降落到地面。

	（a）估算整个地球表面的年平均降雨量（以毫米表示，球面积为 $4\pi R^{2}$）。

	（b）太阳辐射到地球的能量中只有约 $50\%$ 到达地面，$W$ 只是其中的一部分。太阳辐射到地球的能量没能全部到达地面，这是为什么？请说明二个理由。
\end{enumerate}

%% answer:
\jdanswer{
\begin{enumerate}
	\item
	$R=6.37\times10^{6}\Um$
	\item
	（a）$1.0\times10^{3}\Umm$；（b）大气层的吸收，大气层的散射或反射，云层遮挡等。
\end{enumerate}
}
%% memo:
\memoanswer{
（1）由题意知 $2\pi R\times\frac{1}{4}=1.00\times10^{7}$

解得 $R=6.37\times10^{6}\Um$ ①

（2）（a）设太阳在一年中辐射到地球水面部分的总能量为 $W$，有

$W=1.87\times10^{24}\UJ$

凝结成雨滴年降落到地面水的总质量为 $m$

$m=\frac{W\times(1-7\%)\times(1-35\%)}{2.2\times10^{6}}=5.14\times10^{17}\Ukg$ ②

使地球表面覆盖一层水的厚度为 $h$

$h=\frac{m}{\rho S_{\text{地球}}}$

解得 $h=1.01\times10^{3}\Umm$ ③

整个地球表面年平均降雨量约为 $1.0\times10^{3}\Umm$。

（b）大气层的吸收，大气层的散射或反射，云层遮挡等。

评分标准：全题 10 分。第（1）小题 3 分，第（2）小题 7 分。其中（1）得出①给 3 分，写出 $R=6.4\times10^{6}\Um$，同样给分。（2）（a）得出②给 2 分，得出③给 2 分。（b）写出 1 个原因，得 1 分；2 个或 2 个以上正确的原因，得 3 分；如果写出其它合理的原因，也同样给分。
}

\item
%% number: 20
%% paperName: 2001年上海高考第 20 题 10 分
%% typeId: 5
%% type: 辨析题
%% chapter: 运动和力
%% point: 瞬时加速度
%% method: 无
%% score: 10
%% degree: 450
%% duplicateId: 0
%% body:
如图\ref{2001上海20a}所示，一质量为 $m$ 的物体系于长度分别为 $l_{1}$、$l_{2}$ 的两根细线上，$l_{1}$ 的一端悬挂在天花板上，与竖直方向夹角为 $\theta$，$l_{2}$ 水平拉直，物体处于平衡状态。现将 $l_{2}$ 线剪断，求剪断瞬时物体的加速度。
\begin{enumerate}
	\item
	下面是某同学对该题的一种解法：

	解：设 $l_{1}$ 线上拉力为 $T_{1}$，$l_{2}$ 线上拉力为 $T_{2}$，重力为 $mg$，物体在三力作用下保持平衡 $T_{1}\cos\theta=mg$，$T_{1}\sin\theta=T_{2}$，$T_{2}=mg\tan\theta$。剪断线的瞬间，$T_{2}$ 突然消失，物体即在 $T_{2}$ 反方向获得加速度。因为 $mg\tan\theta=ma$，所以加速度 $a=g\tan\theta$，方向在 $T_{2}$ 反方向。你认为这个结果正确吗？请对该解法作出评价并说明理由。
	\item
	若将图\ref{2001上海20a}中的细线 $l_{1}$ 改为长度相同、质量不计的轻弹簧，如图\ref{2001上海20b}所示，其他条件不变，求解的步骤和结果与（1）完全相同，即 $a=g\tan\theta$，你认为这个结果正确吗？请说明理由。
\end{enumerate}

\twopicture[align=right, w=5cm, labelA=2001上海20a, labelB=2001上海20b]
{figs/20a.png}
{figs/20b.png}

%% answer:
\jdanswer{
\begin{enumerate}
	\item
	不正确。因为当 $l_{2}$ 被剪断的瞬间，$l_{1}$ 上的张力 $T_{1}$ 即刻发生了变化，所以 $T_{1}$ 与重力 $mg$ 的合力就不再等于 $mg\tan\theta$，所以加速度 $a\neq g\tan\theta$。
	\item
	正确。因为 $l_{2}$ 被剪断的瞬间，弹簧 $l_{1}$ 的长度未发生变化，$T_{1}$ 大小和方向都不变，小球的合力大小等于 $T_{2}$，方向水平向左，故小球的加速度 $a=g\tan\theta$。
\end{enumerate}
}
%% memo:
\memoanswer{
（1）不正确。因为当 $l_{2}$ 被剪断的瞬间，$l_{1}$ 上的张力 $T_{1}$ 即刻发生了变化，所以 $T_{1}$ 与重力 $mg$ 的合力就不再等于 $mg\tan\theta$，所以加速度 $a\neq g\tan\theta$。

（2）正确。因为 $l_{2}$ 被剪断的瞬间，弹簧 $l_{1}$ 的长度未发生变化，$T_{1}$ 大小和方向都不变，小球的合力大小等于 $T_{2}$，方向水平向左，故小球的加速度 $a=g\tan\theta$。
}

\item
%% number: 21
%% paperName: 2001年上海高考第 21 题 12 分
%% typeId: 6
%% type: 计算题
%% chapter: 气体性质
%% point: 气体多过程
%% method: 临界
%% score: 12
%% degree: 450
%% duplicateId: 0
%% body:
如图所示，一定量气体放在体积为 $V_{0}$ 的容器中，室温为 $T_{0}=300\UK$，有一光滑导热活塞 C（不占体积）将容器分成 A、B 两室，B 室的体积是 A 室的两倍，A 室容器上连接有一 U 形管（U 形管内气体的体积忽略不计），两边水银柱高度差为 $76\Ucm$，右室容器中连接有一阀门 K，可与大气相通。（外界大气压等于 $76\UcmHg$）求：
\begin{enumerate}
	\item
	将阀门 K 打开后，A 室的体积变成多少？
	\item
	打开阀门 K 后将容器内的气体从 $300\UK$ 分别加热到 $400\UK$ 和 $540\UK$，U 形管内两边水银面的高度差各为多少？
\end{enumerate}

\onepicture[align=right, w=6cm]{\includesvg{figs/21}}

%% answer:
\jdanswer{
\begin{enumerate}
	\item
	$V_{A}=\frac{2}{3}V_{0}$
	\item
	$400\UK$ 时，水银柱的高度差为 0；$540\UK$ 时，水银高度差为 $15.2\Ucm$。
\end{enumerate}
}
%% memo:
\memoanswer{
（1）设外界大气压为 $p_{0}=76\UcmHg$，开始时，A 室气体的体积为 $V_{A0}=\frac{1}{3}V_{0}$，由气压平衡知

$p_{A0}=p_{0}+76\UcmHg=2p_{0}$，

打开阀门，A 室气体等温变化，$p_{A}=p_{0}$，体积 $V_{A}$。由玻意耳定律得

$p_{A0}V_{A0}=p_{A}V_{A}$ ①

$V_{A}=\frac{p_{A0}V_{A0}}{p_{A}}=\frac{2}{3}V_{0}$ ②

（2）从 $T_{0}=300\UK$ 升到 $T$，体积为 $V_{0}$，压强为 $p_{A}$，是等压过程。

从 $T_{0}=300\UK$ 升到 $T$ 时活塞恰好运动到容器的最右端，此时体积为 $V_{0}$，压强为 $p_{A1}=p_{0}$，A 室气体做等压变化，由盖-吕克萨定律得

$T=\frac{V_{0}T_{0}}{V_{A}}=\frac{3}{2}\times300\UK=450\UK$ ③

$T_{1}=400\UK<450\UK$，所以从 $T_{0}=300\UK$ 升高到 $T_{1}=400\UK$ 时，A 室气体做等压变化，$p_{A1}=p_{A0}=p_{0}$，水银柱的高度差为 0。

从 $T=450\UK$ 升高到 $T_{2}=540\UK$，A 室气体做等容变化，压强为 $p_{A2}$，由查理定律得

$\frac{p_{A}}{T}=\frac{p_{A2}}{T_{2}}$ ④

$p_{A2}=\frac{T_{2}p_{A}}{T}=1.2p_{0}=91.2\UcmHg$ ⑤

$T_{2}=540\UK$ 时，水银高度差为 $15.2\Ucm$。

评分标准：全题 12 分。第（1）小题 4 分，第（2）小题 8 分。其中（1）得出①、②各得 2 分。（2）得出③式，得 3 分；结果正确，得 3 分。得出④、⑤式，各得 1 分；结果正确，得 2 分。
}

\item
%% number: 22
%% paperName: 2001年上海高考第 22 题 8 分
%% typeId: 6
%% type: 计算题
%% chapter: 电磁感应
%% point: 电磁感应电路分析
%% method: 无
%% score: 8
%% degree: 450
%% duplicateId: 0
%% body:
半径为 $a$ 的圆形区域内有均匀磁场，磁感强度为 $B=0.2\UT$，磁场方向垂直纸面向里，半径为 $b$ 的金属圆环与磁场同心地放置，磁场与环面垂直，其中 $a=0.4\Um$，$b=0.6\Um$，金属环上分别接有灯 $L_{1}$、$L_{2}$，两灯的电阻均为 $R_{0}=2\UO$，一金属棒 MN 与金属环接触良好，棒与环的电阻均忽略不计。
\begin{enumerate}
	\item
	若棒以 $v_{0}=5\Ums$ 的速率在环上向右匀速滑动，求棒滑过圆环直径 $OO'$ 的瞬时（如图所示）MN 中的电动势和流过灯 $L_{1}$ 的电流。
	\item
	撤去中间的金属棒 MN，将右面的半圆环 $OL_{2}O'$ 以 $OO'$ 为轴向上翻转 $90^{\circ}$，若此时磁场随时间均匀变化，其变化率为 $\frac{\Delta B}{\Delta t}=\frac{4}{\pi}$（$\UTs$），求 $L_{1}$ 的功率。
\end{enumerate}

\onepicture[align=right, w=4.8cm]{\includesvg{figs/22}}

%% answer:
\jdanswer{
\begin{enumerate}
	\item
	$E_{1}=0.8\UV$，$I_{1}=0.4\UA$
	\item
	$P_{1}=1.28\times10^{-2}\UW$
\end{enumerate}
}
%% memo:
\memoanswer{
（1）金属棒运动到 $OO'$ 时，切割产生的电动势为

$E_{1}=B\cdot2av=0.2\times0.8\times5\UV=0.8\UV$ ①

流过灯 $L_{1}$ 的电流为

$I_{1}=\frac{E_{1}}{R}=\frac{0.8}{2}\UA=0.4\UA$ ②

（2）由法拉第电磁感应定律得

$E_{2}=\frac{\Delta\Phi}{\Delta t}=0.5\times\pi a^{2}\times\frac{\Delta B}{\Delta t}=0.32\UV$ ③

两灯的电阻相等，故灯 $L_{1}$ 的功率为

$P_{1}=\frac{\left(\frac{E_{2}}{2}\right)^{2}}{R}=1.28\times10^{-2}\UW$ ④

评分标准：全题 13 分。第（1）小题 6 分，第（2）小题 7 分。其中（1）正确得出①式得 3 分，得出②式得 3 分；（2）得出③式 4 分，得出④式得 3 分。
}

\item
%% number: 23
%% paperName: 2001年上海高考第 23 题 18 分
%% typeId: 6
%% type: 计算题
%% chapter: 力物体的平衡
%% point: 力矩平衡
%% method: 无
%% score: 18
%% degree: 350
%% duplicateId: 0
%% body:
如图所示，光滑斜面的底端 a 与一块质量均匀、水平放置的平板光滑相接，平板长为 $2L$，$L=1\Um$，其中心 C 固定在高为 $R$ 的竖直支架上，$R=1\Um$，支架的下端与垂直于纸面的固定转轴 O 连接，因此平板可绕转轴 O 沿顺时针方向翻转。问：
\begin{enumerate}
	\item
	在斜面上离平板高度为 $h_{0}$ 处放置一滑块 A，使其由静止滑下，滑块与平板间的动摩擦因数 $\mu=0.2$，为使平板不翻转，$h_{0}$ 最大为多少？
	\item
	如果斜面上的滑块离平板的高度为 $h_{1}=0.45\Um$，并在 $h_{1}$ 处先后由静止释放两块质量相同的滑块 A、B，时间间隔为 $\Delta t=0.2\Us$，则 B 滑块滑上平板后多少时间，平板恰好翻转。（重力加速度 $g$ 取 $10\Umsq$）
\end{enumerate}

\onepicture[align=right, w=6.5cm]{\includesvg{figs/23}}

%% answer:
\jdanswer{
\begin{enumerate}
	\item
	$h_{0}\le0.16\Um$
	\item
	$t=0.2\Us$
\end{enumerate}
}
%% memo:
\memoanswer{
（1）设 A 滑到 a 处的速度为 $v_{0}=\sqrt{2gh_{0}}$ ①

$f=\mu N$，$N=mg$，$f=ma$，

$a=\mu g$ ②

滑到板上离 a 点的最大距离为 $v_{0}^{2}=2\mu gs_{0}$，

$s_{0}=\frac{2gh_{0}}{2\mu g}=\frac{h_{0}}{\mu}$ ③

A 在板上不翻转应满足条件：摩擦力矩小于正压力力矩，即 $M_{\text{摩擦}}\le M_{\text{压力}}$

$\mu mgR\le mg(L-s_{0})$ ④

$h_{0}\le\mu(L-\mu R)=0.2(1-0.2)=0.16\Um$ ⑤

（2）当 $h=0.45\Um$，$v_{A}=\sqrt{2gh}=\sqrt{2\times10\times0.45}=3\Ums$

$v_{A}=v_{B}=3\Ums$ ⑥

设 B 在平板上运动直到平板翻转的时刻为 $t$，取 $\Delta t=0.2\Us$

$s_{A}=v_{A}(t+\Delta t)-\frac{1}{2}\mu g(t+\Delta t)^{2}$ ⑦

$s_{B}=v_{B}t-\frac{1}{2}\mu gt^{2}$ ⑦$'$

两物体在平板上恰好保持平板不翻转的条件是

$2\mu mgR=mg(L-s_{A})+mg(L-s_{B})$ ⑧

由⑦+⑦$'$式等于⑧式，得 $t=0.2\Us$

评分标准：全题 15 分。第（1）小题 7 分，第（2）小题 8 分。其中（1）得出①、②、③各得 1 分，判断 $M_{\text{摩擦}}\le M_{\text{压力}}$ 正确得 2 分，④、⑤式各得 1 分。（2）得出⑤式得 1 分，①式得 1 分，写出③式得 3 分，最后结果正确得 3 分。
}

\end{enumerate}

\end{document}
